\chapter{סיכום ומפת דרכים}

בניית סטארט-אפ נכון אינה הימור על רעיון בודד אלא תהליך למידה ממושמע תחת
אי-ודאות. ראינו כי רוב הכשלונות נובעים מבעיה לא-אמיתית, ולכן הסדר הנכון הוא
\emph{בעיה לפני פתרון}, \emph{למידה לפני הרחבה}, ו\emph{מזומן לפני יוקרה}.

\section{מפת דרכים תמציתית}

\begin{enumerate}
  \item \textbf{בעיה:} זהה בעיה תכופה וכואבת; אמת אותה בשיחות עם לקוחות אמיתיים.
  \item \textbf{\en{MVP}:} בנה ניסוי מינימלי שבודק את ההשערה המסוכנת ביותר.
  \item \textbf{לולאה:} הפעל \en{Build--Measure--Learn} ומזער את זמן המחזור.
  \item \textbf{\en{PMF}:} חתור להתאמת מוצר-שוק לפני כל השקעה בצמיחה.
  \item \textbf{צוות:} בנה צוות מייסד משלים עם הסכמות ברורות מראש.
  \item \textbf{כלכלה:} ודא ש-$\mathrm{LTV}/\mathrm{CAC} \geq 3$ ונהל את ה-\en{runway}.
  \item \textbf{סיכונים:} בדוק בכל רגע את ההשערה המסוכנת ביותר, בזול ובמהירות.
\end{enumerate}

\section{מילה אחרונה}

הכלים שהוצגו — מתודולוגיית הלקוחות, הלולאה הרזה, כלכלת היחידה וניהול הסיכונים —
אינם נוסחת קסם, אלא מסגרת משמעת. הם מגדילים את הסיכוי שהמשאב היקר ביותר, זמנו
של הצוות, יושקע בבעיה הנכונה. זה, בתמצית, ההבדל בין סטארט-אפ שבונה את המוצר
\emph{נכון} לבין סטארט-אפ שבונה את המוצר \emph{הנכון}.
